\section{Continue from samples to published pixels}
Attach \texttt{drafts/png-packing-construction/attach.py} to your completed row-stage
program. Your loader, collector, inflater and reconstruction instructions are retained.
The attachment supplies final stage-call wiring, an unfinished scalar/MSA packer,
a blank 64-byte mask table and an unfinished result-publication region. Generate a
selected PNG with \texttt{fixtures.py} and select its observation with
\texttt{prepare.py}; these helpers do not replace your algorithm regions.

The actual file caller uses a 512-word destination stride. Direct function examples
also use smaller strides so padding and the final active extent can be examined
without a large display. Check the two geometries separately. The checkpoint scratch
registers below illustrate one cumulative implementation; completed free-writing
checks use the public function contracts and memory effects.

Once pixel stores are being constructed, development status four can accompany
partial or complete displayed pixels. It denotes unfinished work. The completed
format/capacity failure result one must retain the framebuffer, and every final
context result remains unchanged until the last publication region succeeds.
The two diagnostic steps near the end add no code: they inspect the failure and
success boundaries in the supplied caller before final publication is authored.
